\scanpage{49}
\mspage{4}

$\mathbb C$-linéaire, \deleted{\unsure{commutant} d’ailleurs (par tr. de str.)}
\[
(17')\qquad V_{\mathbb C}\xrightarrow{p_\omega}V_{(\omega)}
\]
\margin{cette application commute par transport de str. \unsure{à|:}
$O^+(V)$ (opérant sur $V_{\mathbb C}$ par complexification}

d’où une application \inserted{$\mathbb C$-lin.} canonique
$V_{\mathbb C}\longrightarrow\prod_{\omega\in\underline\omega}V_{(\omega)}$,
qui est un isomorphisme
\[
(18)\qquad p=(p_\omega)_{\omega\in\underline\omega}:
V_{\mathbb C}\xrightarrow{\sim}\prod_{\omega\in\underline\omega}
\underbrace{V_{(\omega)}}_{\simeq V_{\mathbb C}^{\omega}}.
\]
\margin{l’opération d’homothétie $i$ sur $V_{\mathbb C}$
\deleted{\unsure{comme}} \inserted{\unsure{s’identifie}}
\unsure{au prod des} $u_\omega$ sur les facteurs $V_{(\omega)}$}

Les $V_{\omega}$ sont apparus comme
des quotients de $V_{\mathbb C}$, si on veut les
regarder comme sous-espaces\inserted{— \unsure{auquel|auxquels} corresp. la \unsure{not°} $V_{\mathbb C}^{\omega}$ —} on trouve
que
\[
\begin{aligned}
(19)\qquad V_{\mathbb C}^{\omega\deleted{\unsure{\mathbb C}}}
&=\operatorname{Ker}\bigl(V_{\mathbb C}\xrightarrow{p_{\omega'}}V_{(\omega')}\bigr)\\
&=\{x\in V_{\mathbb C}\mid u_{\omega_{\mathbb C}}(x)=i.x\}
\end{aligned}
\]
\margin{\uline{\unsure{Ainsi|Aussi}}
$V_{\mathbb C}^{\omega}=\{\underbrace{x\unsure{-|+}iu_\omega x}_{\substack{\Vert\\\varphi_\omega(x)}}\mid x\in V\}$,
\unsure{où} $\varphi_\omega:V\xrightarrow{\sim}V_{\mathbb C}^{\omega}$
le $\mathbb R$-isom. inverse de $p_\omega$ $(17')$}

En fait, les $V_{\mathbb C}^{\omega}$\inserted{$(\omega\in\underline\omega)$} sont des sous-espaces
propres pour l’action de $O^+(V)$, i.e.
stables pour cette action, tandis que
les éléments de $O^-(V)$ les interchan-
gent.
\margin{vérifier que cette indexation des sous-espaces propres
$V_{\mathbb C}^{\omega}$ coïncide avec celle définie par la conv. de Stokes}

On s’intéresse au fibré vectoriel

\margin{48}
